\begin{derivation}[从\eqnrefs{eq:sm-gim-loop-sum-dp83}到\eqnrefs{eq:sm-gim-mass-suppression-dp83}]
质量本征基中的夸克带电流为
\begin{equation*}
 \Lag_{\rm CC}^{q}
 =-\frac{\hbar c\,g}{\sqrt2}
 \left[
 \bar u_i\gamma^\mu P_LV_{ij}d_j\,\mathcal W_\mu^+
 +\bar d_j\gamma^\mu P_LV_{ij}^*u_i\,\mathcal W_\mu^-
 \right].
\end{equation*}
考虑 $d_j\to d_k$、$j\neq k$ 的一圈味改变中性流。Dyson 展开中任意包含两个带电流顶角的收缩，都含味因子
\begin{equation*}
 V_{ij}V_{ik}^*=V_{ik}^*V_{ij},
\end{equation*}
因为第一个顶角把外部 $d_j$ 连接到内部 $u_i$，第二个厄米共轭顶角再把同一内部味 $u_i$ 连接到外部 $d_k$。

给定外部粒子与动量后，选一组完备、线性无关的 Lorentz--Dirac 算符基 $\{\mathcal K_a(p_{\rm ext})\}$。同一微扰阶的全部规范完成图求和后，任意内部味 $i$ 的贡献都能按这组基展开：
\begin{equation*}
 \mathcal A_{j\to k}^{(i)}
 =V_{ik}^*V_{ij}
 \sum_a\mathcal K_a(p_{\rm ext})F_a(x_i),
 \qquad
 x_i=\frac{m_{u_i}^2}{m_W^2}.
\end{equation*}
这里 $F_a$ 由实际圈积分定义。例如在维数正规化中，每个 $F_a$ 都来自形如
\begin{equation*}
 \mu_{\rm DR}^{2\epsilon}
 \int\frac{\dd^d\ell}{(2\pi)^d}
 \frac{N_a(\ell,p_{\rm ext};m_{u_i},m_W)}
 {\prod_r D_r(\ell,p_{\rm ext};m_{u_i},m_W)}
\end{equation*}
的积分及同阶规范完成项。对内部味求和并交换有限求和与算符基展开，得到
\begin{align*}
 \mathcal A_{j\to k}^{\rm FCNC}
 &=\sum_i\mathcal A_{j\to k}^{(i)}\\
 &=\sum_a\mathcal K_a(p_{\rm ext})
 \left[\sum_iV_{ik}^*V_{ij}F_a(x_i)\right].
\end{align*}
因此
\begin{equation*}
 C_a^{jk}=\sum_iV_{ik}^*V_{ij}F_a(x_i),
\end{equation*}
即正文\eqnrefs{eq:sm-gim-loop-sum-dp83}。

CKM 幺正性 $V^\dagger V=I$ 的非对角元满足
\begin{equation*}
 \sum_iV_{ik}^*V_{ij}=0,
 \qquad j\neq k.
\end{equation*}
对任意参考点 $x_*$，逐个算符系数加减同一个 $F_a(x_*)$：
\begin{align*}
 C_a^{jk}
 &=\sum_iV_{ik}^*V_{ij}
 \left[F_a(x_i)-F_a(x_*)+F_a(x_*)\right]\\
 &=\sum_iV_{ik}^*V_{ij}[F_a(x_i)-F_a(x_*)]
 +F_a(x_*)\sum_iV_{ik}^*V_{ij}\\
 &=\sum_iV_{ik}^*V_{ij}[F_a(x_i)-F_a(x_*)].
\end{align*}
这就是\eqnrefs{eq:sm-gim-subtracted-exact-dp85}。若 $x_i=x_*$ 对所有内部味都成立，则每个独立算符系数 $C_a^{jk}$ 都等于零，因而完整振幅为零。

若 $F_a$ 在连接 $x_*$ 与 $x_i$ 的区间可微，微积分基本定理还给出精确质量差表示
\begin{align*}
 F_a(x_i)-F_a(x_*)
 &=\int_{x_*}^{x_i}\dd y\,F_a'(y)\\
 &=(x_i-x_*)\int_0^1\dd t\,
 F_a'\!\left[x_*+t(x_i-x_*)\right],
\end{align*}
因此
\begin{equation*}
 C_a^{jk}
 =\sum_iV_{ik}^*V_{ij}(x_i-x_*)
 \int_0^1\dd t\,
 F_a'\!\left[x_*+t(x_i-x_*)\right].
\end{equation*}
这条等式直接表明共同的内部质量部分被消去，剩余量由质量劈裂加权。

对轻内部质量取 $x_*=0$。若 $F_a\in C^3([0,x_{\max}])$，Taylor 定理带积分余项给
\begin{equation*}
 F_a(x_i)
 =F_a(0)+x_iF_a'(0)+\frac{x_i^2}{2}F_a''(0)
 +\frac{x_i^3}{2}\int_0^1\dd t\,(1-t)^2F_a^{(3)}(tx_i).
\end{equation*}
代入并使用 $\sum_iV_{ik}^*V_{ij}=0$：
\begin{equation*}
 C_a^{jk}
 =F_a'(0)\sum_iV_{ik}^*V_{ij}x_i
 +\mathcal R_{a,jk}^{(\ge2)},
\end{equation*}
其中
\begin{align*}
 \mathcal R_{a,jk}^{(\ge2)}
 &\equiv
 \frac{F_a''(0)}2
 \sum_iV_{ik}^*V_{ij}x_i^2+R_{a,jk}^{(3)},\\
 R_{a,jk}^{(3)}
 &\equiv
 \frac12\sum_iV_{ik}^*V_{ij}x_i^3
 \int_0^1\dd t\,(1-t)^2F_a^{(3)}(tx_i).
\end{align*}
于是
\begin{align*}
 |R_{a,jk}^{(3)}|
 &\le\frac16
 \max_{0\le y\le x_{\max}}|F_a^{(3)}(y)|
 \sum_i|V_{ik}^*V_{ij}|x_i^3,\\
 |\mathcal R_{a,jk}^{(\ge2)}|
 &\le\frac{|F_a''(0)|}{2}
 \sum_i|V_{ik}^*V_{ij}|x_i^2
 +|R_{a,jk}^{(3)}|.
\end{align*}
最后用 $x_i=m_{u_i}^2/m_W^2$ 即得到正文\eqnrefs{eq:sm-gim-mass-suppression-dp83} 及其余项控制。若 $F_a$ 在零点附近非解析，或某个 $x_i$ 不小，则保留精确差值\eqnrefs{eq:sm-gim-subtracted-exact-dp85}。
\end{derivation}
